% Generated by write_targeted_followups_tex.py; edit the template.
\subsection{H1四函数差的配对伴随分析}
为确定路径超限在推断函数上的表现，将原42项失败与同设备、预算及实际输入的另一MALA执行方式配对，去重后得到80项（42失败、38有效），覆盖13个原输入标签。使用已保存的实际初值、噪声及接受随机数，独立NumPy重放共1130496次转移。它是浮点参照，不是数学精确路径；设备、预算和执行器副本不增加独立重复数。

在原参数尺度分别计算$v/3$、$\tanh(v/3)$、$\ind(x_1>0)$和$\cos(x_1e^{-v/2})$。对全路径及剔除原512步后的保留段，逐链与等长四链合并报告最大绝对差、均方根差、平均绝对差、有符号均值差和事件失配。参数映射复算与保存输出一致；独立接受事件和符号函数失配均为0，3200行函数汇总未出现非有限值。

\begin{table}[htbp]\centering\small
\begin{tabular}{lrrr}\toprule
函数 & 最大逐点绝对差 & 最大合并均值差 & 最大单链均值差\\\midrule
$v/3$ & $1.910\times10^{-3}$ & $1.078\times10^{-7}$ & $4.350\times10^{-7}$\\
$\tanh(v/3)$ & $1.452\times10^{-3}$ & $4.690\times10^{-8}$ & $1.904\times10^{-7}$\\
$\ind(x_1>0)$ & $0$ & $0$ & $0$\\
$\cos(x_1e^{-v/2})$ & $3.269\times10^{-3}$ & $3.078\times10^{-7}$ & $1.232\times10^{-6}$\\
\bottomrule\end{tabular}
\caption{H1保留段在全部80项中的差异极值。均值差列取绝对值；各列最大值未必来自同一任务。合并指同一任务四条等长链，不是把80项混合。完整RMS、平均绝对差、任务身份及全程结果另列于源表。}
\label{tab:h1-function-companion}
\end{table}

连续函数的均值差小于最严重位置的差，提示其影响应按函数和汇总方式评价；平均绝对差与RMS仍用于区分局部性和正负抵消。该伴随分析不改变原数值合格标准、不把失败轨迹计入主结果，也不对80项作独立重复置信区间。冻结分析身份为\texttt{h1-function-path-companion-v1}。全部差数组独立复算后重建3200行，恢复检查新增重放0；完整数据和两幅位置/配对图见仓库的同名分析目录。
